\documentclass{article}
\title{Practice Problem 2.47}
\author{CSSAPP}
\usepackage{booktabs, amssymb, tabularx}

\begin{document}
    \maketitle

    \begin{flushleft}
        Consider a $5-bit$ floating-point representation based on the IEEE 
        floating-point format, with one sign bit, two exponent bits $(k = 2)$, 
        and two fraction bits $(n = 2)$. The exponent bias is $2^{2-1} -1 = 1$.\\
        The table that follows enumerates the entire nonnegative range for this $5-bit$
        floating-point representation. Fill in the blank table entries using the following 
        directions.\\ 
        $e:$        The value represented by considering the exponent field to be an unsigned integer.\\
        $E:$        The value of the exponent after biasing.\\
        $2^E$:      The numeric weight of the exponent.\\
        $f$:        The value of  the fraction\\
        $M$:        The value of the significant.\\
        $2^E x M$:  The (unreduced) fractional value of the number.\\
        $V$:        The reduced fractional value of the number.\\
        Decimal:    The decimal representation of the number.\\
        \smallbreak
        Express the values of $2^E, f, M, 2^E \ast M, V$ either as integers (when possible) or 
        as fractions of the form $\frac{x}{y}$, where $y$ is the power of $2$. \\

        \bigskip
        \begin{tabular}{ccccccccc}
        $Bits$  & $e$   & $E$   & $2^E$ & $f$   & $M$   & $2^E \ast M$ & $V$    & $Decimal$\\
        \midrule
        $0\,00\,00$   &  $0$ & $0$ & $1$ & $\frac{0}{4}$ & $\frac{0}{4}$ & $\frac{0}{4}$ & $0$ & $0.0$\\
        $0\,00\,01$   &  $0$ & $0$ & $1$ & $\frac{1}{4}$ & $\frac{1}{4}$ & $\frac{1}{4}$ & $\frac{1}{4}$ & $0.25$\\
        $0\,00\,10$   &  $0$ & $0$ & $1$ & $\frac{2}{4}$ & $\frac{2}{4}$ & $\frac{2}{4}$ & $\frac{1}{2}$ & $0.5$\\
        $0\,00\,11$   &  $0$ & $0$ & $1$ & $\frac{3}{4}$ & $\frac{3}{4}$ & $\frac{3}{4}$ & $\frac{3}{4}$ & $0.75$\\
        $0\,01\,00$   &  $1$ & $0$ & $1$ & $\frac{0}{4}$ & $\frac{4}{4}$ & $\frac{4}{4}$ & $1$ & $1.0$\\
        $0\,01\,01$   &  $1$ & $0$ & $1$ & $\frac{1}{4}$ & $\frac{5}{4}$ & $\frac{5}{4}$ & $\frac{5}{4}$ & $1.25$\\
        $0\,01\,10$   &  $1$ & $0$ & $1$ & $\frac{2}{4}$ & $\frac{6}{4}$ & $\frac{6}{4}$ & $\frac{3}{2}$ & $1.5$\\
        $0\,01\,11$   &  $1$ & $0$ & $1$ & $\frac{3}{4}$ & $\frac{7}{4}$ & $\frac{7}{4}$ & $\frac{7}{4}$ & $1.75$\\
        $0\,10\,00$   &  $2$ & $1$ & $2$ & $\frac{0}{4}$ & $\frac{4}{4}$ & $2$ & $2$ & $2.0$\\
        $0\,10\,01$   &  $2$ & $1$ & $2$ & $\frac{1}{4}$ & $\frac{5}{4}$ & $\frac{5}{2}$ & $\frac{5}{2}$ & $2.5$\\
        $0\,10\,10$   &  $2$ & $1$ & $2$ & $\frac{2}{4}$ & $\frac{6}{4}$ & $3$ & $3$ & $3.0$\\
        $0\,10\,11$   &  $2$ & $1$ & $2$ & $\frac{3}{4}$ & $\frac{7}{4}$ & $\frac{7}{2}$ & $\frac{7}{2}$ & $3.5$\\
        $0\,11\,00$   &  $3$ & $2$ & $4$ & $\frac{0}{4}$ & $\frac{4}{4}$ & $4$ & $4$ & $4.0$\\
        $0\,11\,01$   &  $3$ & $2$ & $4$ & $\frac{1}{4}$ & $\frac{5}{4}$ & $5$ & $5$ & $5.0$\\
        $0\,11\,10$   &  $3$ & $2$ & $4$ & $\frac{2}{4}$ & $\frac{6}{4}$ & $6$ & $6$ & $6.0$\\
        $0\,11\,11$   &  $3$ & $2$ & $4$ & $\frac{3}{4}$ & $\frac{7}{4}$ & $7$ & $7$ & $7.0$\\
\end{tabular}
    \end{flushleft}
\end{document}